By Lemma~\(\mathrm{lem{:}clw\text{-}population\text{-}design\text{-}rule}\), the source conditions reduce the published population criterion to affine-linear residual variance divided by sample size. Definition~\(\mathrm{def{:}clw\text{-}linear\text{-}cost\text{-}selector}\) maps its support only to \(M_d\) and specializes cost to \(nC_d\), so its continuous criterion is \(C_dr_d^{\mathrm{lin}}(P)/B\).

Since \(X\) has one-point support, \(X_{\mathrm{nr}}\) is zero-dimensional and \(\xi_d=(1,W_x,W_y)^\top\). Symmetry and conditional independence given \(U\) give
\[
 \mathbb E(U)=\mathbb E(W_x)=\mathbb E(W_y)=0,
 \quad \mathbb E(UW_x)=x,
 \quad \mathbb E(UW_y)=y,
 \quad \mathbb E(W_xW_y)=xy.                              \tag{11}
\]
Therefore the full affine Gram, including its unique intercept, is
\[
 \Gamma_d(P_{\chi,\epsilon})
 =\begin{pmatrix}1&0&0\\0&1&xy\\0&xy&1\end{pmatrix},
 \qquad \det\Gamma_d(P_{\chi,\epsilon})=1-x^2y^2>0.       \tag{12}
\]
The outcome cross moment is \((0,x,y)^\top\). Thus the intercept coefficient is zero and inversion of the remaining normal equations gives
\[
 \beta_1=\frac{x(1-y^2)}{1-x^2y^2},
 \qquad
 \beta_2=\frac{y(1-x^2)}{1-x^2y^2}.                        \tag{13}
\]
The projection identity yields
\[
 r_d^{\mathrm{lin}}(P_{\chi,\epsilon})
 =1-(x,y)
 \begin{pmatrix}1&xy\\xy&1\end{pmatrix}^{-1}
 \begin{pmatrix}x\\y\end{pmatrix}
 =\frac{(1-x^2)(1-y^2)}{1-x^2y^2}.                         \tag{14}
\]
Direct enumeration of the four proxy cells gives the same conditional-mean squared error, so (14) also equals \(R_d(P_{\chi,\epsilon})\) at the center. No off-center equality is used. Because \(Y(1)=Y(0)=U\) and \(Z(1)=Z(0)\), the affine residual law is identical in the two arms of the transported randomized comparison; hence arm homoskedasticity holds. Finite support gives the moment condition, and (12) proves \(P_{\chi,\epsilon}\in\mathfrak P_{\mathcal D}^{\mathrm{CLW}}\).

Substitute \((x,y)=(\epsilon,y_\chi)\) for \(d_{\mathrm p}\) and \((x,y)=(1/2,1/2)\) for \(d_{\mathrm e}\). Then
\[
 r_{d_{\mathrm p}}^{\mathrm{lin}}
 =\frac{3}{10\chi}
 \frac{1-\epsilon^2}{1-\epsilon^2(1-3/(10\chi))},
 \qquad r_{d_{\mathrm e}}^{\mathrm{lin}}=\frac35.         \tag{15}
\]
Since \(C_{d_{\mathrm p}}=5\chi\) and \(C_{d_{\mathrm e}}=5\),
\[
 B\Psi_{\mathrm{CLW},d_{\mathrm p}}
 =\frac32\frac{1-\epsilon^2}{1-\epsilon^2y_\chi^2}
 <\frac32<3=B\Psi_{\mathrm{CLW},d_{\mathrm e}}.           \tag{16}
\]
The same substitution proves the unrestricted-risk chain because equality was proved at the center. Hence both continuous selectors choose \(d_{\mathrm p}\).

If \(B\ge2C_d\), the elementary floor inequality gives
\[
 0\le \frac{B}{\lfloor B/C_d\rfloor}-C_d
 \le\frac{2C_d^2}{B}.                                     \tag{17}
\]
Multiplying by \(r_d^{\mathrm{lin}}(P)\) shows
\[
 B\Psi_{\mathrm{CLW},d}^{\mathrm{int}}(P;B)
 \longrightarrow C_dr_d^{\mathrm{lin}}(P).                \tag{18}
\]
Finiteness of \(\mathcal D\) and the strict gap in (16) furnish the asserted finite \(B_0(P_{\chi,\epsilon})\).

All center cells are strictly positive. Consistency is built into the potentials; experimental treatment is randomized; observational treatment is latent-exchangeable with probabilities in \([3/8,5/8]\); proxy kernels are source-invariant; and conditional proxy independence implies the sequential restriction. A binary proxy with nonzero signal has a nonsingular conditional-expectation matrix. Since \(\epsilon,y_\chi,1/2\ne0\), both completeness conditions and both bridge operators hold. The displayed finite bridges solve their equations, so \(P_{\chi,\epsilon}\in\mathfrak P_{\mathcal D}\). Since \(Y(1)=Y(0)=U\), \(\tau=0\), and both bridge means vanish.

The bridge formulas give
\[
 V_{E,d}=\frac4{y^2},
 \qquad
 V_{O,d}=\frac{64(1-y^2)(1+14x^2)}{225x^2y^2}.             \tag{19}
\]
Consequently
\[
 V_{E,d_{\mathrm p}}=\frac{40\chi}{10\chi-3},
 \qquad
 V_{O,d_{\mathrm p}}
 =\frac{64}{75(10\chi-3)}(\epsilon^{-2}+14),              \tag{20}
\]
whereas \(V_{E,d_{\mathrm e}}=16\) and \(V_{O,d_{\mathrm e}}=384/25\). The budget-profile theorem yields
\[
 \mathcal V_{d_{\mathrm e}}^*
 =\left(\sqrt{2\cdot16}+\sqrt{3\cdot384/25}\right)^2
 =\frac{3872}{25}.                                         \tag{21}
\]
Substitution gives the exact ratio (8) in the proposed statement. Since \(p_{E,1}=p_{O,1}=1/2\),
\[
 V_{g,d}^{\mathrm{IKM}}=\frac14V_{g,d},
 \qquad g\in\{E,O\},                                     \tag{22}
\]
so every IKM profile is one quarter of its arm-normalized counterpart and the factor cancels from the ratios. The named and generic center selectors coincide, proving all four equalities.

Dropping the first nonnegative square-root term gives
\[
 \rho\ge\frac{5\chi+1}{363(10\chi-3)}(\epsilon^{-2}+14)
 >\frac{\epsilon^{-2}+14}{726},                            \tag{23}
\]
because the difference between the two coefficients is \(5/\{726(10\chi-3)\}>0\). Put
\[
 A(\chi)=10\chi+1+\frac3{10\chi-3},
 \quad D(\chi)=\frac{16}{75}+\frac{16}{15(10\chi-3)},
 \quad C_0=\frac{25}{3872}.
\]
The exact ratio satisfies
\[
 \epsilon^2\rho
 =C_0\left[\epsilon^2A(\chi)+D(\chi)(1+14\epsilon^2)
 +2\epsilon\sqrt{1+14\epsilon^2}\sqrt{A(\chi)D(\chi)}\right]. \tag{24}
\]
For \(\chi=k\epsilon^{-2/3}\), locally uniformly on compact subsets of \((0,\infty)\),
\[
 \epsilon^2\rho
 =\frac1{726}+C_0\epsilon^{2/3}
 \left\{2\sqrt{\frac{32k}{15}}+\frac8{75k}\right\}
 +o(\epsilon^{2/3}).                                       \tag{25}
\]
The bracket diverges at both endpoints, so rescaled minimizers are compact. Its derivative changes sign once, at \(k_\star=\sqrt[3]{18}/15\), and \(C_0q(k_\star)=5\sqrt[3]{12}/968\). This proves the optimizer scale and sharp expansion.